\documentclass[12pt]{article}
\usepackage[utf8]{inputenc}
\usepackage[T1]{fontenc}
\usepackage{hyperref}
\title{Wadjigu Language Resources for K-12 Curriculum Development}
\author{Compiled by Claude Sonnet 5 (claude-sonnet-5) for Andr\'as Kornai}
\date{}
\begin{document}
\maketitle

\section{Introduction}
Wadjigu (ISO 639-3 wdu) is an extinct Aboriginal Australian language of inland eastern Queensland, around the Expedition Range and Woorabinda. Wikipedia lists it as a dialect in the Bidjara group, and notes that linguists do not agree on all of the Bidjara dialects. Very little was found. No wordlist, grammar or teaching material turned up in the search.

\section{Primary Educational Resources}
None found.

\section{Linguistic and Reference Materials}
\begin{itemize}
\item Wikipedia, Wadjigu people (background only): \url{https://en.wikipedia.org/wiki/Wadjigu_people}
\item Wikipedia, Bidjara language (the group the language is assigned to): \url{https://en.wikipedia.org/wiki/Bidjara_language}
\item AIATSIS AUSTLANG record E39: \url{https://collection.aiatsis.gov.au/austlang/language/E39}
\end{itemize}

\section{Religious and Community Resources}
None found.

\section{Digital and Audio Resources}
None found. The State Library of Queensland open word lists did not appear to include Wadjigu: \url{https://data.gov.au/data/dataset/3fa7a159-f280-48f9-9a83-9209f73942ef}

\section{Curriculum Development Context}
The University of Queensland Fryer Library holds handwritten anthropological field notes catalogued under Wadjigu people (E39), including material by Caroline Tennant-Kelly. The library warns that they contain sensitive content: \url{https://manuscripts.library.uq.edu.au/index.php/wadjigu-people-e39}. Any curriculum would have to start from archival material and from the Bidjara-related community, and should be developed with the relevant Traditional Owners.
\end{document}